% 02-related-work.tex — truth: 本文件（2026-09-10 交割：R22 后全部正文修订直改 tex，原装配源 md 已冻结为历史设计稿，见其头部治理注记）
% 纪律: 三小节主题式 + 每小节 gap 句收尾; tab1 三近邻差距表随行;
% [CITATION-NEEDED] 内嵌标记移除, 作者待补状态记录 refs.bib note(3 条: yolo2025petx/mct2024tip/gra2024tnnls)。
% R3-1/R2-2 修订 (2026-11-15): 新增 §2.4 逐条划界 TREBA/BehaveMAE/ProSL + tab:compmap 失配表;
%   三处 Gap 改写承认动物姿态自监督已有工作; 首次性主张收窄为机制组合级。
% 2026-11-16 回修: §2.4 新增与 class-incremental/continual learning 的定位段（审查#5）,
%   refs.bib 补 rebuffi2017icarl / wang2024continual / masana2023cil 三条经核验题录。

\section{Related work}\label{sec:related}

Animal behavior recognition supports applications ranging from livestock welfare monitoring to working-dog training. Training a single service or assistance dog costs tens of thousands of dollars, and more than half of candidates fail to graduate~\citep{grimm2026science}, which motivates objective, automated behavior assessment. Deployed systems already combine wearable sensors with classical machine learning for canine posture estimation~\citep{plosone2023imu}, yet the field lacks a recognition framework that works when labels are scarce and evaluation criteria keep evolving. We review three bodies of work that bear on this gap---animal behavior understanding, skeleton-based self-supervised learning, and pseudo-label-driven semi-supervised learning---and then position our contribution against the three nearest neighbors and against class-incremental learning (Section~\ref{ssec:rw-position}).

\subsection{Animal behavior understanding}\label{ssec:rw-animal}

Existing animal behavior pipelines divide into appearance-based, pose-based, and sensor-based lines.

Appearance-based methods classify behavior directly from RGB frames or video. YOLO-PetX adapts a detection backbone to anomalous dog behavior~\citep{yolo2025petx}. TP-CanineNet applies temporal contrastive learning with pseudo-labels to canine separation-anxiety video~\citep{tpcanine2025}. These methods inherit the strengths of image recognition but depend on large annotated corpora, and their accuracy degrades under background change, occlusion by handlers, and breed variation in appearance.

Pose-based methods first extract keypoints, then classify dynamics. DeepLabCut~\citep{mathis2018dlc} and SLEAP~\citep{pereira2022sleap} made markerless animal pose estimation practical. Building on such inputs, ASBAR unifies skeleton-based animal action recognition with PoseConv3D and reports 75.3\% top-1 on primate benchmarks~\citep{asbar2024}. BCST-GCN adapts ST-GCN-style graph convolutions with bidirectional cross-attention to pig behavior and reaches 95.36\% average accuracy on its four-behavior dataset compiled from online public videos~\citep{bcstgcn2026}. Hierarchical representation learning of dog behavior from single-view 3D pose has also been explored~\citep{miyai2025workshop}. Skeleton representations are compact, background-invariant, and robust across individuals---properties that should suit low-resource settings---but every behavior classifier above is trained fully supervised on a fixed label set.

Sensor-based methods use accelerometers or IMUs. They capture coarse daily activities reliably~\citep{plosone2023imu} but cannot resolve fine-grained behaviors that differ mainly in body configuration rather than global motion.

\paragraph{Public data landscape.}
The scarcity that motivates this work is visible in the data that actually exists. Among the datasets we could identify, the only large-scale real animal-behavior video benchmark with per-video action labels is Animal Kingdom~\citep{ng2022animalkingdom}, in which dogs appear as a sparsely-labeled subset; canine-specific resources are otherwise wearable inertial time series with ethogram-annotated behaviors~\citep{vehkaoja2022dogsensor,plosone2023imu}, synthetic dog videos built for pose estimation~\citep{shooter2024sydogvideo}, or pose-only still images. To our knowledge none provides real working-dog behavior as skeleton sequences with per-segment expert action labels at benchmark scale. By ``scarce'' we mean skeleton-registered expert behavior labels: Animal Kingdom provides frame-level action labels but no 3D skeletons, so obtaining skeleton sequences requires self-extraction---it is the joint availability of clean skeletons and per-segment action labels, not temporal labels per se, that is missing. Our public-real tier must weakly exploit exactly this, which is why it is assembled by self-extraction from Animal Kingdom (Section~\ref{sec:datasets}) and why the target-domain real-K9 tier is pre-registered yet left unexercised. On our own evidence (Section~\ref{sec:ablation}; the alignment analysis in Section~\ref{sec:conclusion}, L9) the negative end-to-end boundary on the canine tier is attributable to label misalignment and extraction quality in the only public source of canine behavior video with action labels; whether it also reflects true domain-level scarcity is precisely the question the pre-registered K9 pilot is designed to separate.

\paragraph{Gap.}
Current animal pipelines assume a fixed taxonomy and supervision-heavy training. Animal pose self-supervision is no longer unexplored: TREBA learns rodent behavior representations from a fraction of expert labels by programming pretext tasks from domain knowledge~\citep{sun2021treba}, and h/BehaveMAE extends masked autoencoders to spatio-temporal motion-capture sequences, exposing a movement-to-activity hierarchy~\citep{stoffl2024behavemae}. Domain-adaptive pretraining improves primate behavior recognition~\citep{primateda2025}, and self-supervised video pretraining extends to multi-animal settings~\citep{mabe22,beast2025}. What these lines share is that they stop at representations and their evaluation probes: none couples them to temporal behavior proposals, to a label-expansion loop, or to a mechanism for absorbing taxonomy change---no surveyed system is designed so that changing evaluation criteria can be handled without re-annotating the corpus.

\subsection{Skeleton-based self-supervised and unsupervised learning}\label{ssec:rw-ssl}

Self-supervised representation learning for skeletons is mature in the human domain. ST-GCN established graph convolutions over joint topologies~\citep{yan2018stgcn}, and PoseConv3D recasts keypoints as heatmap volumes for 3D convolution~\citep{duan2022posec3d}. Among purely unsupervised objectives, AimCLR maximizes asymmetry across extreme augmented views with an InfoNCE~\citep{vandenoord2018representation} queue~\citep{guo2022aimclr}. Under the linear-evaluation protocol (frozen encoder + FC-softmax classifier; three-stream joint/motion/bone fusion), the conference paper reports 78.9\% on NTU60 cross-subject~\citep{guo2022aimclr}, and the official repository's re-measurement of the released checkpoints gives 79.18\%. Its journal extension is credited with 80.9\% on NTU60 cross-subject linear evaluation in the official repository's results table~\citep{aimclrpp2024}. Cross-view contrastive objectives form a parallel self-supervised family for skeletons~\citep{li2021crossclr}, and we treat the encoder choice as replaceable within our framework; AimCLR serves as the reference physics-layer instantiation rather than a load-bearing design commitment. For temporal structure, unsupervised segmentation methods quantize motion into discrete vocabularies; a motion-word quantization pipeline (SMQ) derives behavior-proposal boundaries from quantized motion-word sequences~\citep{gokay2025smq}.

Two observations motivate our treatment of this literature. First, virtually all benchmarks, augmentations, and graph templates assume the human body plan; transferring these objectives to quadruped kinematics requires adapting joint topology, augmentation priors, and view construction. Second, pretraining and segmentation are studied separately, although a recognition system needs both: representations without proposals cannot localize behaviors, and proposals without semantics cannot name them. State-space backbones such as VideoMamba~\citep{videomamba2024} and Mamba-MSQNet for animal action recognition~\citep{fazzari2024mambamsqnet} offer efficiency alternatives but do not address label scarcity either.

\paragraph{Gap.}
Self-supervised skeleton machinery is mature for humans and, since TREBA and h/BehaveMAE~\citep{sun2021treba,stoffl2024behavemae}, has entered the animal domain---but as representation learning, validated through probes and clustering, and still decoupled from unsupervised temporal decomposition and from any label-expansion loop; our repository-scale survey found no published animal-skeleton pipeline combining self-supervised representation learning with the latter two. In the human skeleton domain, masked autoencoders~\citep{skeletonmae2022} and foundation-model efforts~\citep{skeletonfm2025} pretrain representations for broad downstream transfer; our decoupling claim is orthogonal---the frozen component is the physics encoder, the revisable component is the taxonomy-facing semantic layer, and the measured quantity is taxonomy-transition cost rather than representation quality (systematic repository-scale survey; full matrix in Appendix~\ref{app:novelty}).

\subsection{Semi-supervised, few-shot, and pseudo-label learning}\label{ssec:rw-semi}

Semi-supervised skeleton action recognition in the human domain relies mostly on consistency regularization and teacher--student schemes, including momentum-contrastive teachers~\citep{mct2024tip} and graph-representation alignment~\citep{gra2024tnnls}. These methods regularize predictions under perturbation but receive no external guidance about which classes matter. MAC-Learning introduces anchor-based contrastive learning, where ``anchors'' are contrastive sample pairs selected to structure the embedding space~\citep{mac2022learning}; it does not maintain seed-annotated anchors that supervise cluster assignment, and it includes no iterative pseudo-label loop. Closest in mechanism, ProSL iterates between clustering and pseudo-label generation in self-supervised contrastive learning: starting from a contrastively pretrained encoder, a codebook built by clustering supplies cluster-level pseudo-labels that are progressively refined over rounds~\citep{qin2025prosl}. ProSL targets fully unsupervised representation learning on human benchmarks, whereas PSD runs a \emph{semi-supervised recognition} loop driven by externally supplied rule-engine seed anchors---the semantic layer, not the representation, is the object being trained---and adds confidence filtering, temporal behavior proposals, and the taxonomy-transition cost framing. Point-supervised skeleton action segmentation is another close neighbor: it generates prototype-similarity pseudo-labels with clustering from one labeled frame per segment~\citep{pointsup2026}; our combination differs in rule-engine seed anchors, confidence-filtered iteration, and the taxonomy-transition cost framing---we claim firstness for that combination, not for prototype pseudo-labels per se. Few-shot approaches such as PAINet learn episode-based metric classifiers~\citep{painet2023} and likewise lack a mechanism for expanding labels beyond the support set.

Vision foundation models supply another route to low-shot recognition~\citep{caron2021dino}. Skeleton-to-image encoding maps skeletons into images consumed by VFM-pretrained encoders~\citep{skel2img2026}, demonstrating that foundation-model features transfer to skeleton data. This line uses VFMs for representation only; it does not close the loop from a handful of annotated seeds to full taxonomy coverage. In the animal domain specifically, TP-CanineNet is the closest use of pseudo-labels, operating on RGB video rather than skeletal dynamics~\citep{tpcanine2025}.

\paragraph{Positioning against the nearest neighbors.}
Table~\ref{tab:neighbors} summarizes the nearest neighbors of our semantic-layer design.

\begin{table}[pos=t]
\centering
\caption{Nearest neighbors of the semantic-layer design and their mechanism coverage. None of the surveyed rows couples seed-annotated anchors, cluster-based pseudo-label iteration, and a self-supervised physics encoder for temporal-skeleton recognition.}
\label{tab:neighbors}
\small
\setlength{\tabcolsep}{4pt}
\begin{tabularx}{\linewidth}{p{2.9cm}p{2.0cm}ccc p{1.9cm}}
\hline
Work & Domain & Seed anchors & Cluster + pseudo-iter & SSL enc. & Output \\
\hline
MAC-Learning~\citep{mac2022learning} & Human skel. & $\times$ (pairs) & $\times$ & $\times$ & Repr. \\
ProSL~\citep{qin2025prosl} & Human skel. & $\times$ (clustering) & \checkmark\ (unsuper.) & \checkmark & Repr. \\
Skeleton-to-Image~\citep{skel2img2026} & Skel.\ repr. & $\times$ & $\times$ & \checkmark & Repr. \\
TP-CanineNet~\citep{tpcanine2025} & Canine RGB & $\times$ & Partial & $\times$ & Cls. \\
\textbf{Ours} & \textbf{Quadruped stream} & \checkmark\ (rule-eng.) & \checkmark\ (conf.-filt.) & \checkmark\ (physics) & Recog.+seg. \\
\hline
\end{tabularx}
\end{table}

To the best of our knowledge, transferring the combination of seed-annotated anchor learning, prototype clustering, and iterated confidence-filtered pseudo-labeling to temporal-skeleton recognition has not been reported. Our claim of firstness rests on a repository-scale survey covering ten query groups and eighteen candidate works---twenty-two across the final matrix's eleven rows once the submission-time re-verification row and the revision's three nearest-neighbor rows are included---with zero occupancy (full matrix in Appendix~\ref{app:novelty}), complemented by an arXiv/CrossRef re-verification at submission time that surfaced point-supervised skeleton segmentation~\citep{pointsup2026} as the nearest neighbor---differentiated above and in Section~\ref{ssec:rw-position}. We state this claim for the \emph{mechanism combination}; we do not claim firstness for any individual ingredient (Section~\ref{ssec:rw-position}).

\subsection{Positioning against the three nearest neighbors}\label{ssec:rw-position}

The three nearest neighbors of this work are TREBA~\citep{sun2021treba}, h/BehaveMAE~\citep{stoffl2024behavemae}, and ProSL~\citep{qin2025prosl}. Table~\ref{tab:compmap} maps them, with two further neighbors, against the data and protocol they were built for, and names the specific mismatch that prevents a direct comparison.

\begin{table}[pos=t]
\centering
\caption{Comparability map: nearest neighbors, their data and protocol, and the specific mismatch that prevents a head-to-head comparison with PSD. ``Superv.\ type'' is the supervision signal the method's main result is reported under. None of the rows is directly comparable to our three-tier protocol without re-experimentation, which we therefore do not claim as a comparison.}
\label{tab:compmap}
\small
\setlength{\tabcolsep}{4pt}
\hyphenpenalty=10000
\emergencystretch=1.5em
\begin{tabularx}{\linewidth}{p{2.65cm}p{2.3cm}p{2.3cm}p{1.9cm}X}
\hline
Method & Data / domain & Task \& protocol & Superv.\ type & Protocol mismatch vs.\ ours \\
\hline
TREBA~\citep{sun2021treba} & Mouse fly-courtship / rodent home-cage \textbf{RGB video + pose} & Representation learning; $k$-NN / classification probes, label-efficiency curves & Self-sup. + few expert labels & Image/pose features; no temporal proposals, no semantic layer, no taxonomy-change setting \\
h/BehaveMAE~\citep{stoffl2024behavemae} & Animal mocap (mice, macaques, ASTM) \textbf{pose seqs.} & Unsup.\ masked pretraining; linear probes at three levels & Unsupervised & Pretraining quality only; no recognition pipeline, no pseudo-label loop, no transition-cost view \\
ProSL~\citep{qin2025prosl} & Human skeletons (NTU60/120, Kinetics) & Unsup.\ contrastive learning with iterated cluster pseudo-labels; linear eval & Unsupervised (cluster pseudo-labels) & Trains the \emph{representation}; no seed anchors, no confidence-filtered recognition loop, human-only \\
MAC-Learning~\citep{mac2022learning} & Human skeletons (NTU, Kinetics) & Semi-sup.\ representation learning via anchor contrastive pairs & Partial labels & Anchors = contrastive pairs, not rule-derived seeds; no cluster pseudo-label iteration \\
Point-supervised seg.~\citep{pointsup2026} & Human skeletons & Action segmentation from one labeled frame per segment & Point labels & Single-frame labels, not rule seeds; segmentation only; no frozen-physics decoupling \\
\textbf{Ours (PSD)} & \textbf{Quadruped streams + public human/ape benchmarks} & \textbf{Three-tier pre-registered protocol; recognition + segmentation under taxonomy transitions} & \textbf{Rule-engine seeds (tens of clips)} & --- \\
\hline
\end{tabularx}
\end{table}

\paragraph{TREBA~\citep{sun2021treba} (CVPR 2021).}
\emph{What it does.} Task programming injects domain knowledge into self-supervised pretext tasks (frame--behavior affinity, temporal-order prediction) so that rodent behavior representations can be learned from limited expert annotations; the learned representations are evaluated with $k$-NN and light classifiers, and task programming approaches expert-level labeling efficiency on fly-courtship and mouse home-cage videos. \emph{Where we differ.} TREBA learns an embedding for RGB/pose frames under a fixed ethogram; it has no temporal behavior-proposal stage, no mechanism that converts embeddings into taxonomy-wide labels, and no notion of an evaluation taxonomy that later changes. \emph{PSD increment.} PSD uses its label-free physics layer \emph{as a substrate}: the semantic layer turns embeddings plus unsupervised proposals into labels under the current taxonomy, and the measured quantity is the cost of changing that taxonomy, not probe accuracy.

\paragraph{h/BehaveMAE~\citep{stoffl2024behavemae} (ECCV 2024).}
\emph{What it does.} A hierarchical masked autoencoder for animal motion-capture data (mice, macaques, and the multi-species ASTM benchmark) learns spatio-temporal behavior representations without labels; probes show lower layers encode fine-grained movements and higher layers complex activities. \emph{Where we differ.} h/BehaveMAE is a representation-quality study: its outputs are latents and probe accuracies, not a recognition pipeline. It does not segment continuous streams into behavior proposals, does not expand labels, and does not consider taxonomy evolution. \emph{PSD increment.} PSD treats such encoders as the replaceable frozen physics layer (Section~\ref{ssec:rw-ssl}) and contributes the missing half: a revisable semantic layer and an explicit account of what happens---to accuracy and to wall-clock cost---when the taxonomy changes.

\paragraph{ProSL~\citep{qin2025prosl} (Machine Learning 2025).}
\emph{What it does.} ProSL improves unsupervised skeleton-based action recognition by iterating between clustering and pseudo-label generation: a contrastively pretrained encoder is clustered into a codebook, cluster assignments serve as pseudo-labels, and this loop progressively refines the representation; it reports the strongest published unsupervised results on human skeleton benchmarks. \emph{Where we differ.} ProSL's pseudo-labels are cluster identities used to train the \emph{encoder}; there are no externally supplied seeds, so the label space is the clustering's own---and nothing connects it to a domain taxonomy, to behavior proposals from unsupervised segmentation, or to a transition-cost question. \emph{PSD increment.} PSD inverts the direction of trust: representations are frozen after pretraining, and pseudo-labels grow a \emph{semantic layer} from rule-engine seed anchors toward a prescribed taxonomy under confidence filtering; the evaluated quantities are low-budget recognition accuracy and taxonomy-transition cost. This is also why ProSL cannot be transplanted as a baseline into our three-tier protocol without re-experimentation (Table~\ref{tab:compmap}).

\paragraph{Relation to class-incremental and continual learning.}
Taxonomy evolution under a frozen backbone is adjacent to class-incremental learning (CIL) and continual learning, where a model absorbs new classes over a sequence of tasks without catastrophically forgetting old ones~\citep{rebuffi2017icarl,wang2024continual,masana2023cil}. The two settings share machinery---frozen or slowly updated representations, head re-training, and logit-level reconciliation across label spaces (classifier realignment and logit aggregation for merged heads)---but their objectives differ. CIL methods are designed and benchmarked for incremental accuracy under sequences of disjoint class sets, typically spending rehearsal buffers, distillation, or parameter isolation to resist forgetting~\citep{rebuffi2017icarl,masana2023cil}; PSD asks a different question: the \emph{cost--benefit accounting} of an explicitly declared taxonomy remap $\mathcal{Y}\to\mathcal{Y}'$ over an already-collected unlabeled stream, where the semantic layer rebuilds coverage from rule-engine seeds plus unsupervised proposals and no rehearsal or memory-preservation objective is in play. CIL baselines therefore answer a question our protocol does not ask---accuracy retention across incremental sessions under an arriving-class stream---and are not comparable within our matched-budget, frozen-remap transition protocol without re-experimentation; conversely, our transition-cost endpoints (wall-clock and annotation units at matched accuracy under a frozen remap) are not quantities CIL benchmarks report. We position the two lines as complementary: CIL supplies mechanisms for absorbing \emph{arriving} classes under stability--plasticity trade-offs, while this paper measures when the decoupled structure makes an \emph{arbitrary declared remap}---merges and renames here; splits and additions measured on the synthetic tier (Section~\ref{sec:experiments})---worth absorbing at low cost.

\paragraph{Gap.}
Three absences compound across the three lines reviewed above. Animal behavior understanding offers no label-efficient path and no mechanism for taxonomy change (Section~\ref{ssec:rw-animal}). Skeleton-based self-supervision supplies representation machinery that has begun to reach animals (TREBA, h/BehaveMAE) but stops at representations, and treats segmentation as a separate problem (Section~\ref{ssec:rw-ssl}). Semi-supervised learning supplies pseudo-label machinery---most closely ProSL's cluster-iteration---whose components have, to our knowledge, never been assembled for temporal-skeleton recognition under externally supplied seed anchors (Section~\ref{ssec:rw-semi}). No surveyed method couples a label-free physics layer (self-supervised dynamics modeling plus unsupervised segmentation) with an anchor-guided semantic layer that grows from tens of seeds to a complete, revisable taxonomy. This coupling is the design whose cost--benefit boundary this paper measures.
